\section{总结与延伸}

\subsection{核心论点回顾}

整场66分钟的对话中，Karpathy 反复回到几个核心论点：

\begin{enumerate}
    \item \textbf{工作范式已经不可逆地改变}——从2024年12月起，软件工程师的默认工作流完全不同。衡量生产力的单位从代码行变为 Token 吞吐量。
    \item \textbf{把自己从循环中移除}——AutoResearch 不是一个项目，而是一种哲学。有客观指标的地方，人类不应该是瓶颈。
    \item \textbf{AI 的参差性是结构性的}——在 RL 优化范围内是超级智能，范围外几乎停滞。这不会因为模型变大而自动解决。
    \item \textbf{数字空间优先，物理世界滞后}——翻转比特比加速物质快一百万倍，这决定了变革的节奏。
    \item \textbf{教育从教人变为教 Agent}——你的价值在于那几个 Agent 产生不了的关键 bit。
\end{enumerate}

\subsection{三个值得持续关注的信号}

\begin{itemize}
    \item \textbf{AutoResearch 的分布式版本（Open Ground）能否真正运转}——这将直接影响前沿研究是继续集中在少数实验室，还是分散到互联网 Agent 集群。
    \item \textbf{模型的参差性是暂时的还是结构性的}——如果在可验证领域的进步确实不能泛化到其他领域，那自主 Agent 的可靠性天花板会比预期更低。
    \item \textbf{前沿实验室进一步中心化的趋势}——当决策权集中在越来越少的人手里，Karpathy 所说的"不能说的话"会变成什么？
\end{itemize}

\subsection{拓展阅读}

\begin{itemize}
    \item AutoResearch 开源仓库：\url{https://github.com/karpathy/auto-research}（2026年3月开源）
    \item MicroGPT 项目：Karpathy 的200行极简 GPT 实现（2026年2月发布）
    \item OpenClaw：Peter Steinberger 的开源自主 Agent 框架
    \item Daniel Suarez,《Daemon》：Karpathy 推荐的关于超级智能与社会重组的科幻小说
    \item No Priors 播客完整视频：\url{https://www.youtube.com/watch?v=kwSVtQ7dziU}
\end{itemize}

%% ==================== 正文结束 ==================== %%

\end{document}
